% titel: Blatt 6, A4: Wann gilt die Oktettregel nicht?
% skript: Mortimer Kap. 8 (chemie: Atombindung, Oktettregel)

Die Oktettregel ist eine Faustregel für Hauptgruppen-Nichtmetalle. Elemente der
2.\ Periode ($\mathrm{C,N,O,F}$) erfüllen sie strikt; Elemente ab Periode~3
können ihr Oktett erweitern, Hauptgruppenelemente der 2.\ Periode auch
Elektronenmangel (Elektronendefizit) aufweisen.

\begin{werkzeugkasten}[Prüfschema Oktettregel]
\begin{itemize}
  \item Zentralatom der 2.\ Periode ($\mathrm{C,N,O,F}$): nie $>8$
        Valenzelektronen.
  \item Zentralatom ab Periode~3 (P, S, Cl, $\dots$): Oktetterweiterung $>8$
        möglich, z.\ B. $\mathrm{PCl_5}$, $\mathrm{SF_6}$.
  \item Elektronenmangel: $\mathrm{Be}$ und $\mathrm{B}$ bleiben oft bei
        $<8$ (z.\ B. $\mathrm{BF_3}$).
  \item H und He: Oktett gilt nicht, Duett statt dessen ($\mathrm{H_2O}$ ist
        trotzdem eine gültige gewöhnliche Verbindung, nur eben kein
        Oktett-Problem).
\end{itemize}
\textit{Skript:} Mortimer Kap. 8 (Oktettregel).
\end{werkzeugkasten}

\begin{achtung}[Darauf achten]
\begin{itemize}
  \item Eine Verbindung kann die Oktettregel \textit{verletzen}, ohne
        instabil zu sein (z.\ B. $\mathrm{PCl_5}$ und $\mathrm{BF_3}$ sind
        stabile Substanzen).
  \item Entscheidend ist die Zahl der Valenzelektronen am \textit{Zentralatom}
        nach der Lewis-Struktur, nicht die Zahl der Atome.
  \item $\mathrm{H_2O}$ und $\mathrm{CH_4}$ erfüllen die Regel --- sie sind
        Beispiele, bei denen alle Atome ein gültiges Oktett (H-Duett) haben.
\end{itemize}
\end{achtung}
